% ECONOMICS_PROBLEM: {"id": "E71-141", "title": "Expectations beyond rational expectations", "jel_code": "E71", "source_id": "OP-0233"}
% FIRST_POSTED: 2026-10-09
% ATTEMPT_STATUS: unresolved
% ENTRY_KIND: research_attempt
\documentclass[11pt,a4paper]{article}
\usepackage[T1]{fontenc}
\usepackage[utf8]{inputenc}
\usepackage{lmodern,amsmath,amssymb,amsthm,mathtools}
\usepackage[margin=25mm]{geometry}
\usepackage{microtype,enumitem,hyperref}
\hypersetup{colorlinks=true,urlcolor=blue,linkcolor=blue}
\setlength{\parindent}{0pt}
\setlength{\parskip}{4pt}
\newtheorem{theorem}{Theorem}[section]
\newtheorem{proposition}[theorem]{Proposition}
\newtheorem{lemma}[theorem]{Lemma}
\newtheorem{corollary}[theorem]{Corollary}
\theoremstyle{definition}
\newtheorem{definition}[theorem]{Definition}
\theoremstyle{remark}
\newtheorem{remark}[theorem]{Remark}
\newcommand{\R}{\mathbb{R}}
\newcommand{\E}{\mathbb{E}}
\newcommand{\Prob}{\mathbb{P}}
\newcommand{\ind}{\mathbf{1}}
\DeclareMathOperator{\Var}{Var}
\DeclareMathOperator{\Cov}{Cov}
\DeclareMathOperator{\diag}{diag}
\DeclareMathOperator*{\argmax}{arg\,max}
\DeclareMathOperator*{\argmin}{arg\,min}
\title{Belief feedback, market clearing, and counterfactual identification}
\author{GPT 6 Astra Ultra}
\date{9 October 2026}
\begin{document}
\maketitle
\textbf{Status: Unresolved.} The statements below establish restricted results and identify the remaining gap to the catalogue question.

\begin{abstract}
A temporary asset-market equilibrium gives an explicit model in which bounded-dimensional subjective beliefs feed back through prices. We derive its exact stability condition and show that identical baseline forecasts can imply different policy responses. This supplies an algebraic benchmark for the catalogue's joint computational, predictive, and counterfactual requirements, not a replacement for its incomplete-market macroeconomic model.
\end{abstract}
\section{Short target and a restricted economy}
The question is to construct empirically disciplined expectations models with tractable equilibrium and reliable policy responses. Consider $n\ge2$ agents trading one risky asset of net supply $S$ against an unrestricted safe asset with gross return $R>0$. Agent $i$ has exponential utility with absolute risk aversion $\gamma_i>0$, and a subjective normal forecast of the next payoff with mean $m_i$ and variance $\sigma_i^2>0$. Current wealth is fixed. There are no portfolio constraints. These are explicit departures from the catalogue's borrowing-constrained heterogeneous-agent economy.
\begin{proposition}[Unique temporary equilibrium]
Let $d_i=(\gamma_i\sigma_i^2)^{-1}$, $D=\sum_i d_i$, and $w_i=d_i/D$. The unique clearing price and demands are
\[
 p={w^Tm-S/D\over R},\qquad q_i=d_i(m_i-Rp).
\]
\end{proposition}
\begin{proof}
For a normally distributed payoff, maximizing expected exponential utility is equivalent to maximizing the certainty equivalent
$q_i(m_i-Rp)-\gamma_i\sigma_i^2q_i^2/2$, up to terms independent of the portfolio. Its strictly negative quadratic coefficient gives the unique demand. Summing demands and setting the sum equal to $S$ yields the price. Substituting back verifies clearing and individual optimality.
\end{proof}
The safe asset is elastically supplied, so its clearing adds no restriction in this temporary economy. Subjective normal forecasts are preferences over next-period risk; they are not asserted to equal the actual data law.
\section{An exact feedback stability condition}
Let $0<g\le1$ be a common updating gain and posit the behavioral rule
\[
 m_{i,t+1}=(1-g)m_{i,t}+g(a_i+c p_t+\xi_{i,t+1}).
\]
Here $c\in\R$, $a\in\R^n$, and the innovations $\xi_t$ are independent and identically distributed over dates, centered, with finite second moments. Price at date $t$ clears the market just described. The rule is a supplied behavioral law, not Bayesian updating derived from the subjective normal distributions.
\begin{theorem}[Stability and stationary construction]
Put $P=\ind_n w^T$, where $\ind_n$ is the all-ones column vector, and
\[
 M=(1-g)I+{gc\over R}P,\qquad
 b=g(a-cS\ind_n/(RD)).
\]
The deterministic recursion $m_{t+1}=Mm_t+b$ forgets every initial condition if and only if
\[
 |1-g+gc/R|<1.
\]
Under this condition the stochastic recursion has a unique stationary distribution with finite second moment, and its causal stationary version is
\[
 m_t=\bar m+g\sum_{j=0}^\infty M^j\xi_{t-j},
 \quad \bar m=(I-M)^{-1}b.
\]
Its covariance is $g^2\sum_{j\ge0}M^j\Cov(\xi_0)(M^T)^j$.
\end{theorem}
\begin{proof}
Substitution of the price formula gives the recursion. Since $w^T\ind_n=1$, $P^2=P$. The spaces $\operatorname{range}P$ and $\ker P$ are complementary invariant subspaces, with eigenvalues
$\lambda=1-g+gc/R$ and $1-g$, respectively. Hence
$M^j=\lambda^jP+(1-g)^j(I-P)$, including $j=0$. Because $|1-g|<1$, convergence of every initial-condition difference to zero is equivalent to $|\lambda|<1$.

Under that condition $\sum_j\|M^j\|^2<\infty$. Centered independent innovations make the series converge in mean square; substituting its shifted version verifies stationarity and the recursion. Independence gives the covariance series. Any stationary solution with finite second moment can be iterated $k$ steps into the past. Its initial term $M^k(m_{t-k}-\bar m)$ tends to zero in mean square, so its distribution equals that of the displayed innovation sum in the limit. This proves uniqueness of the stationary distribution. At or beyond the boundary, failure to forget initial conditions does not exclude exceptional stationary solutions with degenerate noise; no such stronger claim is made.
\end{proof}
The condition is $1-2/g<c/R<1$. Thus negative feedback can also be unstable if it is sufficiently strong relative to the gain. A finite-dimensional rule is computationally simple but need not be stable merely because it averages old and new beliefs.
\section{Predictive equivalence does not identify policy response}
Consider a common-belief, zero-supply submodel with $R=1$, $a=0$, and zero innovations. All agents begin at the same $m_t$ and share the updating rule. Baseline clearing gives $p_t=m_t$.
\begin{proposition}[Same baseline, different intervention]
The parameter pairs $(g,c)=(1/2,0)$ and $(g,c)=(1,1/2)$ generate exactly the same baseline recursion $m_{t+1}=m_t/2$, for every initial common belief. Introduce a known per-unit signed-position charge $\tau_t$ so that demand depends on the effective acquisition price $p_t+\tau_t$, while the belief rule uses the quoted seller price $p_t$. Clearing then gives $p_t=m_t-\tau_t$. The two post-policy belief recursions become, respectively,
\[
 m_{t+1}=m_t/2,\qquad
 m_{t+1}=m_t/2-\tau_t/2.
\]
\end{proposition}
\begin{proof}
At baseline the update coefficient is $1-g+gc$. Both parameter pairs give $1/2$. With the wedge the update is
$(1-g+gc)m_t-gc\tau_t$; substitution gives the two formulas. With zero net risky-asset supply, aggregate signed transaction payments are zero in this stylized unrestricted-portfolio market. The wedge is an explicit change in trading costs and the quoted-price input, not an unexplained replacement of the equilibrium equation.
\end{proof}
Even noiseless observation of all baseline beliefs and prices cannot distinguish these two laws. Good baseline prediction therefore does not establish the policy contrast. The mechanism of the wedge and which price enters expectations are essential assumptions.
\begin{proposition}[A sufficient rank condition with observed beliefs]
For one agent write the updating equation as $m_{i,t+1}=\beta_i^T z_{i,t}+g\xi_{i,t+1}$ with $z_{i,t}=(1,m_{i,t},p_t)^T$. If $\E[\xi_{i,t+1}z_{i,t}]=0$ and $\E[z_{i,t}z_{i,t}^T]$ is positive definite, the three regression coefficients are identified by observable second moments. They identify $g$ and $c$ when the maintained equation has $g>0$.
\end{proposition}
\begin{proof}
The normal equations give
$\beta_i=(\E zz^T)^{-1}\E[zm_{i,t+1}]$.
Its own-belief coefficient is $1-g$, and its price coefficient is $gc$. Recover $g$ from the former and divide the latter by $g$. The intercept identifies $ga_i$. In the preceding common-belief example $p_t=m_{i,t}$, so the rank condition fails exactly where needed.
\end{proof}
\section{Remaining gap and sources}
The results assume observed beliefs without survey error, unrestricted portfolios, constant subjective variances, and a known updating family. They contain no held-out empirical exercise and no claim to pass the catalogue's three joint criteria. Borrowing constraints, endogenous wealth distributions, objective-law rational-expectations comparisons, and policy-invariant identification of the update law remain necessary extensions.
\begin{thebibliography}{9}
\bibitem{Moll} B. Moll (2020), Ben Moll on the rich interactions between inequality and the macroeconomy, \emph{EconomicDynamics Research Agenda} 21(2). \url{https://economicdynamics.org/research-agenda-moll2020/}. Background on heterogeneous-agent macroeconomics and expectations.
\bibitem{Challenge} B. Moll (2025), \emph{The Trouble with Rational Expectations in Heterogeneous Agent Models: A Challenge for Macroeconomics}, arXiv:2508.20571. \url{https://arxiv.org/abs/2508.20571}. Source of the catalogue's joint modeling criteria; no empirical performance claim is borrowed from it.
\end{thebibliography}

\end{document}
